% Source-derived chapter 06: Proof of \texorpdfstring{\cref{Thm:main}}{main theorem}
% Original source: pixtons-formula-abel-jacobi-picard-stack
\chapter{Proof of \texorpdfstring{\cref{Thm:main}}{main theorem}}
\label{pixtons-formula-abel-jacobi-picard-stack-chapter-06}
\source{pixtons-formula-abel-jacobi-picard-stack}

\begin{remark}[Source section]
\label{pixtons-formula-abel-jacobi-picard-stack-chapter-06-source}
\source{pixtons-formula-abel-jacobi-picard-stack}
\source{pixtons-formula-abel-jacobi-picard-stack}
This chapter reproduces the corresponding section of the retrieved
LaTeX source, including its definitions, statements, examples, and
proofs. It has not yet been translated into Lean.
\notready
\end{remark}

% The source section title was: Proof of \texorpdfstring{\cref{Thm:main}}{main theorem}
\label{sec:proof_of_main_theorem}
\source{pixtons-formula-abel-jacobi-picard-stack}
\subsection{Overview} \label{sec:proof_of_main_theorem_overview}
\source{pixtons-formula-abel-jacobi-picard-stack}
We prove here the main result of the paper:
for $A=(a_1,\ldots,a_n)\in \mathbb{Z}^n$  satisfying $$\sum_{i=1}^n a_i=d\, ,$$
the universal twisted double ramification cycle is calculated by Pixton's formula
$$\mathsf{DR}^{\mathsf{op}}_{g,A} =\P_{g,A,d}^g\
\in {\mathsf{CH}}^g_{\mathsf{op}}(\Picabs_{g,n,d})\, .$$



The result is an equality in the operational Chow group, and therefore
an equality on every finite type family of prestable curves. Given $C\to B$ a prestable curve and a line bundle $\ca L$ on $C$ of relative degree $d$, we obtain
a map $$\phi_{\ca L}\colon B \to \Picabs_{g,n,d}\, .$$
We must prove 
\begin{equation}\label{eq:main_conj}
\source{pixtons-formula-abel-jacobi-picard-stack}
\DRop_{g,A}(\phi_{\ca L})= {\P_{g,A,d}^g}(\phi_{\ca L}) : \Chow_*(B) \to \Chow_{*-g}(B)\, .
\end{equation} 



As explained in \cref{sec:intro_deg_0},  the result for general $A \in \mathbb{Z}^n$ can be reduced  to the case $n=0, d=0$, though the case of arbitrary $A$ will be important in the proof as we proceed through a sequence of special cases. We recall that this reduction used the invariances \invnrn{} and \invnrA{} for the double ramification cycle and Pixton's formula. Note that these will be proved separately and independent of \cref{Thm:main} in \cref{Sect:proofInvariance}, so no circular reasoning occurs.  




\subsection{On an open subset of $\Mbar_{g,n}(\mathbb{P}^l,\beta)$}\label{sec:projective_target}
\source{pixtons-formula-abel-jacobi-picard-stack}
%for \texorpdfstring{$X=\mathbb{P}^l$}{X=PP\^{}l} and  \texorpdfstring{$\mathcal{L}=\ca O(1)$}{L=O(1)}}\label{sec:projective_target}


\newcommand{\kk}{l}
As before, let $A=(a_1,\ldots,a_n) \in \bb Z^n$ with $$\sum_{i=1}^n a_i=d\,.$$ 
We consider here the target $X = \bb P^\kk$.
Let $\beta$ be the class of $d$ times a line in $\bb P^\kk$. 
Let 
$$\ca C \to \Mbar_{g,n}(\mathbb{P}^l,\beta)$$
be the universal curve over the moduli of stable maps to $\mathbb{P}^l$, let 
$$f\colon \ca C \to \mathbb{P}^l$$ be the universal map, and 
let $\ca L = f^*\ca O_X(1)$.


We have a tautological map
\begin{equation}
\phi_{\ca L} \colon \Mbar_{g,n}(\mathbb{P}^l,\beta) \to \Picabs_{g,n,d}. 
\end{equation}
We would like to prove an equality of operational classes 
$$\phi_{\ca L}^*\DRop_{g,A}=\phi_{\ca L}^*{\P_{g,A,d}^g}\, \in \mathsf{CH}_{\mathsf{op}}^g(\Mbar_{g,n}(\mathbb{P}^l,\beta))\, .$$
We will apply  the main result of \cite{Janda2018Double-ramifica} which
relates the double ramification cycle there to 
Pixton's formula. However, 
only the action of $\phi_{\ca L}^*{\P_{g,A,d}^g}$ on the
virtual fundamental class
$[ \Mbar_{g,n}(\mathbb{P}^l,\beta)]^\textup{vir}$ is computed in \cite{Janda2018Double-ramifica}. 
Since we are interested here in the full operational class $\phi_{\ca L}^*{\P_{g,A,d}^g}$,
our first idea is to restrict to the open locus 
$$\Mbar_{g,n}(\mathbb{P}^l,\beta)' \hra \Mbar_{g,n}(\mathbb{P}^l,\beta)$$ 
where (on each geometric fibre) we have $H^1(C, \ca L) = 0$. 

\begin{lemma}
\source{pixtons-formula-abel-jacobi-picard-stack}
\notready\label{lem:fun_eq_vfc}
\source{pixtons-formula-abel-jacobi-picard-stack}
On the smooth Deligne-Mumford stack $\Mbar_{g,n}(\mathbb{P}^l,\beta)'$, the fundamental and virtual fundamental classes coincide. 
\end{lemma}
\begin{proof}
It suffices to show that $H^1(C, f^*T_{\bb P^\kk}) = 0$ on $\Mbar_{g,n}(\mathbb{P}^l,\beta)'$. Pulling back the Euler exact sequence on $\bb P^\kk$ via $f$ yields 
\begin{equation}
0 \to \ca O_C \to \oplus_{1}^{\kk+1} f^*\ca O_{\bb P^\kk}(1) \to f^*T_{\bb P^\kk} \to 0\, .
\end{equation}

 Taking cohomology yields the exact sequence
\begin{equation}
\oplus_{1}^{\kk+1} H^1(C, f^*\ca O_{\bb P^\kk}(1)) \to H^1(C, f^*T_{\bb P^\kk}) \to H^2(C, \ca O_C)\, .
\end{equation}
But $H^1(C, f^*\ca O_{\bb P^\kk}(1)) = 0$ by assumption, and $H^2(C, \ca O_C) = 0$ for dimension reasons. %, so $H^1(C, f^*T_{\bb P^\kk}) = 0$ as required. 
\end{proof}

The next Lemma depends on a careful comparison of the logarithmic and rubber approaches to double ramification cycles, which will be postponed to Section \ref{sec:comparing_stacks}. 
\begin{lemma}
\source{pixtons-formula-abel-jacobi-picard-stack}
\notready\label{cor:operational_equality_projective}\label{lem:projective_space_comparison}
\source{pixtons-formula-abel-jacobi-picard-stack}
Let ${\phi'}_{\ca L}$ be the restriction of ${\phi}_{\ca L}$ to $\Mbar_{g,n}(\mathbb{P}^l,\beta)'$. We have an equality of operational classes 
\begin{equation} \label{eqn:phiprimepullbackDR}
\source{pixtons-formula-abel-jacobi-picard-stack}
{\phi'}_{\ca L}^*\DRop_{g,A}=
{\phi'}_{\ca L}^*{\P_{g,A,d}^g}  \, \in \, \CHop^g(\Mbar_{g,n}(\mathbb{P}^l,\beta)')\, .
\end{equation}
\end{lemma}
\begin{proof}

By \cref{lem:op_eq_ch_for_smooth}, the two sides of  \eqref{eqn:phiprimepullbackDR} are equal if and only if their actions on the fundamental class $[\Mbar_{g,n}(\mathbb{P}^l,\beta)']$ are equal in $\Chow_*(\Mbar_{g,n}(\mathbb{P}^l,\beta)')$. 
By  \eqref{eqn:pixton_vs_uni_pixton}, the action of the right side of \eqref{eqn:phiprimepullbackDR} on 
$$[\Mbar_{g,n}(\mathbb{P}^l,\beta)']=[\Mbar_{g,n}(\mathbb{P}^l,\beta)']^\textup{vir}$$ 
equals the restriction of $\DR_{g,A}(\mathbb{P}^l,\ca L)$ to 
$\Mbar_{g,n}(\mathbb{P}^l,\beta)'$.

The cycle $\DR_{g,A}(\mathbb{P}^l,\ca L)$
is defined in \cite{Janda2018Double-ramifica} as the pushforward of the virtual fundamental class of the space of rubber maps{\footnote{Rubber maps will
be discussed in \cref{sec:prestable_rubber_maps}.}}.  
By \cref{prop:vfc_comparison_summary} of 
\cref{sec:comparing_virtual_classes}, the restriction of
$\DR_{g,A}(X,\ca L)$ to $\Mbar_{g,n}(\mathbb{P}^l,\beta)'$ is equal to ${\phi'}_{\ca L}^*\DRop_{g,A}( [\Mbar_{g,n}(\mathbb{P}^l,\beta)'])$. 

\end{proof}





\subsection{For sufficiently positive line bundles}\label{sec:positive_line_bundle}
\source{pixtons-formula-abel-jacobi-picard-stack}

Let $\pi\colon C\to B$ be an $n$-pointed prestable curve over a scheme of finite type over ${\field}$. Let
$\ca L$ on $C$ be a line bundle 
of relative degree $d$. 
Let $$A = (a_1,\ldots, a_n) \in\mathbb{Z}^n$$ with $\sum_{i=1}^n a_i = d$.  
The line bundle $\ca L$ induces a map 
$$\phi_{\ca L} \colon B \to \Picabs_{g,n,d}\, .$$

We say $\ca L$ is \emph{relatively sufficiently positive} if $\ca L$ is relatively base-point free  and satisfies $R^1\pi_*\ca L = 0$. 


\begin{lemma}
\source{pixtons-formula-abel-jacobi-picard-stack}
\notready\label{sec:very_ample}
\source{pixtons-formula-abel-jacobi-picard-stack}
Let $\ca L$ be a line bundle which is relatively sufficiently positive. 
Then we have an equality  
\begin{equation}\label{eq:very_ample}
\source{pixtons-formula-abel-jacobi-picard-stack}
\DRop_{g,A}(\phi_{\ca L})=
\P_{g,A,d}^g(\phi_{\ca L})  \colon \Chow_*(B) \to \Chow_{*-g}(B)\, .
\end{equation}

\end{lemma}
\begin{proof}
%By \cref{lem:check_equality_on_integral} we may and do assume that $B$ is integral. 

For any finite-type scheme $B$ the union of irreducible components of $B$ maps properly and surjectively to $B$. Thus the pushforward from the Chow groups of the irreducible components to that of $B$ is surjective, and hence it suffices to show the equality \eqref{eq:very_ample} of maps of Chow groups for $B$ irreducible. 

By relative sufficient positivity, 
$$R\pi_*\ca L = \pi_*\ca L$$ 
is a vector bundle on $B$ of rank $N$. 
For a positive integer $l$, we define
\begin{equation}\label{ppp12}
\source{pixtons-formula-abel-jacobi-picard-stack}
E_l =\bigoplus_1^{l+1} R\pi_*\ca L \, , 
\end{equation}
a vector bundle on $B$ of rank $r=N(l+1)$. 
Let $U_l \sub E_l$ denote the open locus of linear systems which are base-point free. Via pullback along $\psi\colon U_l \to B$, we obtain
 a map $$\psi^*\colon \mathsf{CH}_*(B) \to \mathsf{CH}_{*+r}(U_l)\, .$$ 


We claim that for $l>\dim B$, the pullback \eqref{ppp12}
is injective. 
To prove the injectivity, we show that the boundary $E_l \setminus U_l$ has codimension in $E_l$ greater than $\dim B$. Since $E_l \to B$ is flat with irreducible target, it suffices to bound the codimension  on each geometric fibre over $B$: for a  prestable curve $C/{\field}$ and a sufficiently
positive line bundle $\ca L$ on $C$, we must show that the locus in $\bigoplus_1^{l+1} H^0(C, \ca L)$ consisting of base point free linear systems has 
a complement of codimension 
greater than $\dim B$. 

Since $\ca L$ is base point free on $C$, the dimension of the locus in
$\bigoplus_1^{l+1} H^0(C, \ca L)$ where the linear system has
a base point at some given $p\in C$ is $(N-1)(l+1)$. Hence, as $p$ varies, the complement of the
base point free locus in $\bigoplus_1^ {l+1} H^0(C, \ca L)$ has dimension at
most $1 + (N-1)(l+1)$. So the codimension is at least
$$ N(l+1) - 1 - (N-1)(l+1)  = l \, .$$


We have a canonical map $g\colon C \times_B U_l \to \bb P^l$ with 
$g^*\ca O_{\bb P^l}(1) = \ca L$ 
which induces a map $$U_l \to \Mbar_{g,n}(\bb P^l,\beta)$$ which factors via the locus $$\Mbar_{g,n}(\bb P^l, \beta)' \subset \Mbar_{g,n}(\bb P^l, d)$$
where $H^1(C, f^*\ca O_{\bb P^l}(1)) = 0$. 
By construction, the composition $$U_l \xrightarrow{\psi} B \xrightarrow{\phi_{\ca L}} \Picabs_{g,n,d}$$ then factors through the map $ \Mbar_{g,n}(\bb P^l, \beta)' \to \Picabs_{g,n,d}$ induced by the line bundle $f^*\ca O_{\mathbb{P}^l}(1)$ as before. In other words, we have a commutative diagram
%\jocomment{Might the following diagram be helpful to see what is going on:}
\[
\begin{tikzcd}
U_l \arrow[r] \arrow[d,"\psi"]& \Mbar_{g,n}(\bb P^l, \beta)' \arrow[d,"\phi_{f^*\ca O_{\mathbb{P}^l}(1)}"]\\
B \arrow[r,"\phi_{\ca L}"] & \Picabs_{g,n,d}\, . 
\end{tikzcd}
\]


\Cref{cor:operational_equality_projective} then implies that 
\begin{equation}
\left(\DRop_{g,A}-\P_{g,A,d}^g \right)(\phi_{\ca L}\circ \psi)\colon \Chow_*(U_l) \to \Chow_{* - g}(U_l)
\end{equation}
is the zero map, and we conclude the proof of the Lemma from the commutative diagram
\begin{equation}
 \begin{tikzcd}
  \Chow_{*+r}(U) \arrow[rrrr, "\left(\DRop_{g,A}-\P_{g,A,d}^g \right)(\phi_{\ca L}\circ \psi)"]&& && \Chow_{* + r-g}(U)\\
  \Chow_*(B) \arrow[rrrr, "\left(\DRop_{g,A}-\P_{g,A,d}^g\right)(\phi_{\ca L})"] \arrow[u,"\psi^*",hookrightarrow]& &&&\Chow_{* - g}(B) \arrow[u,"\psi^*",hookrightarrow]\, . 
\end{tikzcd}
\end{equation}

\end{proof}



\subsection{With sufficiently many sections}\label{sec:many_markings}
\source{pixtons-formula-abel-jacobi-picard-stack}

Let $\pi\colon C\to B$ be an $n$-pointed prestable curve 
with markings $p_1,\ldots, p_n$
over a scheme of finite type over ${\field}$. Let
$\ca L$ on $C$ be a line bundle 
of relative degree $d$. 
Let $$A = (a_1,\ldots, a_n) \in\mathbb{Z}^n$$ with $\sum_{i=1}^n a_i = d$.  
The line bundle $\ca L$ induces a map 
$$\phi_{\ca L} \colon B \to \Picabs_{g,n,d}\, .$$



\begin{lemma}
\source{pixtons-formula-abel-jacobi-picard-stack}
\notready\label{lem:many_markings}
\source{pixtons-formula-abel-jacobi-picard-stack}
For every geometric fibre of $C/B$, suppose the
complement of the union of irreducible components which carry markings
is a disjoint union of trees of nonsingular
rational curves on which $\ca L$ is trivial. Then we have an equality 
\begin{equation}\DRop_{g,A}(\phi_{\ca L})=
\P_{g,A,d}^g(\phi_{\ca L}) \colon \Chow_*(B) \to \Chow_{*-g}(B)\, .
\end{equation}
%$\phi_{\ca L}^*{Pixton}^{g, r=0} = \phi_{\ca L}^*\DRop$ 
\end{lemma}

\begin{proof}
We can choose $A' = (a'_1,\dots, a'_n)$ with entries 
$$a'_i \gg 0\, , \ \ \ \sum_{i=1}^na_i' =d'\, $$
large enough so that $$\ca L' = \ca L\Big(\sum_{i=1}^n a'_ip_i\Big)$$
is relatively sufficiently positive (by Riemann-Roch for singular curves). 

We obtain an associated map 
$$\phi_{{\ca L}'} \colon B \to \Picabs_{g,n,d+d'}\, .$$
By \cref{sec:very_ample},
\begin{equation} \label{eqn:DRPequalityL42} \DRop_{g,A+A'}(\phi_{{\ca L}'})=
\source{pixtons-formula-abel-jacobi-picard-stack}
\P_{g,A+A',d+d'}^g(\phi_{{\ca L}'}) \colon \Chow_*(B) \to \Chow_{*-g}(B).
\end{equation}
% Since $\phi_{\ca L'} = \tau_{A'} \circ \phi_{\ca L}$ for the map
% $$
% \tau_{A'}\colon \Picabs_{g,n,d} \to \Picabs_{g,n,d+d'} 
% $$
% defined by twisting by $\ca O_C(\sum_{i=1}^n a_i' p_i)$ introduced in \cref{sec:varying_n_and_A}, we have by \cref{Lem:varying_A} and \cref{wwwsss}
Invariance \invnrA{} of \cref{Sect:introinvariance} (proven
in \cref{Sect:proofInvariance}) implies 
\begin{align*}
\DRop_{g,A+A'}(\phi_{{\ca L'}})
%= \tau_{A'}^*\DRop_{g,A+A'} (\phi_{{\ca L}})
&=\DRop_{g,A}(\phi_{{\ca L}})\, ,\\ \P^g_{g,A+A',d+d'}(\phi_{{\ca L'}})
%= \tau_{A'}^*\P^g_{g,A+A',d+d'}(\phi_{{\ca L}})
&=\P^g_{g,A,d}(\phi_{{\ca L}})\, ,     
\end{align*} 
which together with \eqref{eqn:DRPequalityL42} finishes the proof.
\end{proof}

\subsection{Proof in the general case}\label{sec:general_case}
\source{pixtons-formula-abel-jacobi-picard-stack}

To conclude the proof of \cref{Thm:main},
will use the invariances  of \cref{Sect:introinvariance}
(proven in \cref{Sect:proofInvariance}). As discussed in \cref{sec:proof_of_main_theorem_overview}, we can reduce to showing the result in the case $n=0, d=0$.\footnote{In genus $g=1$, we  follow a slightly modified strategy since there we must avoid the case $n=0$ for technical reasons (see Remark \ref{Rmk:genus1n0}). Instead, we can use the invariances to reduce to the case $g=1$, $n=1$, and $A=(0)$. All the proofs below generalize in a straightforward way since the vector $A=(0)$ does not affect the line bundles involved.}



Let $B$ be an irreducible scheme of finite type over ${\field}$.
Let $\pi\colon C\to B$ a prestable curve, and let
$\ca L$ on $C$ be a line bundle 
of relative degree $0$. 
The line bundle $\ca L$ induces a map 
$$\phi_{\ca L} \colon B \to \Picabs_{g,0,0}\, .$$

 \begin{lemma}
\source{pixtons-formula-abel-jacobi-picard-stack}
\notready There exists an \label{ALTt}
\source{pixtons-formula-abel-jacobi-picard-stack}
 alteration{\footnote{An alteration here is a proper, surjective, generically finite morphism between irreducible schemes.}}
 $B' \to B$ and a destabilisation
 \begin{equation}\label{kk3355}
\source{pixtons-formula-abel-jacobi-picard-stack}
 C' \to C \times_B B'
 \end{equation}
 such that $C'$ admits sections $p_1, \ldots, p_m$ and satisfies the
 following property:
 $$(C'/B', p_1, \ldots, p_m)$$ is a family of $m$-pointed prestable curves 
 and 
 for every geometric fibre of $C'/B'$,  the
complement of the union of irreducible components which carry markings
is a disjoint union of trees of nonsingular
rational curves which are contracted by the morphism \eqref{kk3355}.

\end{lemma}

\begin{proof}
We first claim, after an alteration %\{This alteration is actually a modification, i.e. birational. Is it clearer to say this? }\jocomment{} 
$\widehat{B}\to B$,
there exists a multisection\footnote{By a multisection of $C_{\widehat{B}}\to \widehat{B}$,
we mean a closed substack $Z \subset C_{\widehat{B}}$ such that $Z \to \widehat{B}$ is finite and flat.} 
$$Z \subset \, C_{\widehat{B}}=C\times_B \widehat{B}\,  \to \widehat{B}$$ satisfying the following
two conditions:
\begin{enumerate}[label=\roman*)]
    \item[(i)] Over the generic point of $\widehat{B}$, $Z$ is contained in the smooth locus of $C_{\widehat{B}} \to \widehat{B}$.
    \item[(ii)] Every component of every geometric fibre of $C_{\widehat{B}}\to \widehat{B}$  carries at least two  distinct \'etale multisection points in the smooth locus. In other words, the \'etale locus of $Z \to \widehat{B}$ meets the smooth locus of every component of every geometric fibre of $C_{\widehat{B}} \to \widehat{B}$ in at least two points.
\end{enumerate}
To prove the above claim, we observe that for every geometric point $b$ of $B$ there exists an \'etale map $U_p \to B$ and a factorisation $U_p \to C$ whose image meets
the smooth locus of every irreducible component of every geometric fibre in some Zariski neighbourhood $V_b\sub B$ of $b$ at least twice. Choose a finite set of $b$ such that the $V_b$ cover $B$, define $U$ to be the union of the $U_b$, and define $Z'$ to be the closure of the image of $U$ in $C$.  Then $Z' \to B$ is proper and generically finite. 
Let $$\widehat B \to B$$ 
be a modification which flattens $Z'$ (see \cite{Raynaud-Gruson}), and let $Z$ be the strict transform of $Z'$ over $\widehat B$. Then $$Z \to \widehat B$$
is proper, flat, and generically finite, and hence  finite -- so condition (i) is
satisfied. Moreover, $U$ already satisfies condition (ii), and the strict transform of a flat map is just the fibre product, hence $Z$ also satisfies condition (ii). 

Let $\widetilde B \to \widehat B$ be an alteration such that over $\widetilde B$ the multisection $Z$ becomes a disjoint union of sections. 
In other words the pullback 
$$C_{\widetilde{B}}= C \times_B \widetilde B \to \widetilde B$$
has sections $\widetilde p_1, \ldots, \widetilde p_m$ such that, as a set, 
the preimage of $Z$ is given by the union of the images of sections $\widetilde p_1, \ldots, \widetilde p_m$. Such a $\widetilde B$ exists{\footnote{The base $\widehat B$ is excellent since it is finite type over a field.}}
by \cite[Lemma 5.6]{Jong1996Smoothness-semi}.
We can assume that the sections $\widetilde p_i$ are pairwise disjoint over the generic point of $\widetilde B$.

By assumption (i) above, the family $C_{\widetilde B} \to \widetilde B$ with sections $\widetilde p_1, \ldots, \widetilde p_m$ is generically a stable $m$-pointed curve (since 
every component has at least \emph{two} of the sections). We therefore obtain
a rational map $$\widetilde B \dashrightarrow \Mbar_{g,m}\, .$$
Let $B' \to \widetilde B$ be a blow-up resolving the indeterminacy of this map\footnote{As usual, this blowup is constructed by taking the closure of the graph and flattening. Then we check that this ensures the existence of the map to $C_{B'}$ as written below.}

\begin{equation}
\begin{tikzcd}
B' \arrow[d] \arrow[r] & \Mbar_{g,m} \\
\widetilde B \arrow[ru, dashed]&
\end{tikzcd}    
\end{equation}
 and let $C' \to B'$ with sections $p_1, \ldots, p_m : B' \to C'$
 be the pullback of the universal curve over $\Mbar_{g,n}$ to $B'$. Let 
 $$C_{B'} = C \times_B B'$$ 
 be the pullback of $C/B$ under $B' \to \widetilde B \to \widehat B \to B$. Then we have a map $f: C' \to C_{B'}$ fitting in a commutative diagram
 \begin{equation} \label{eqn:partstabilizationalteration}
\source{pixtons-formula-abel-jacobi-picard-stack}
 \begin{tikzcd}
 C' \arrow[rr,"f"] \arrow[rd] & & C_{B'} \arrow[ld]\\
 & B' \arrow[ul,"p_i", bend left] \arrow[ur, "\widetilde p_i", bend right, swap] &
 \end{tikzcd}
 \end{equation}
 such that $f$ is a partial destabilization. On geometric fibers of $C'\to B'$,
 $f$ collapses trees of rational curves to either nodes or coincident
  sections $\widetilde p_i$ on
 %\jocomment{propose changing "markings of" to "sections $\widetilde p_i$ on" since the $\widetilde p_i$ are not really markings} 
 the geometric fibers of $C_{B'}$.
 
 To conclude, we must show that for every geometric point $b \in B'$ and every irreducible component $D \subset C'_b$ which is not contracted by $f$, we can 
 find a marking 
 $$p_i(b)\in D \, .$$ 
 The image of $D$ under $f$ is a component of $(C_{B'})_b$.
 By condition (ii) above, $f(D)$
 has at least one $\widetilde p_i(b)$ in the smooth locus of $f(D)$
 pairwise distinct from all other $\widetilde p_j(b)$.
 Since there are no components of $C'_b$ which collapse to $\widetilde p_i(b)$, we
 must have $p_i(b) \in D$.
 \end{proof}
 


\begin{lemma}
\source{pixtons-formula-abel-jacobi-picard-stack}
\notready \label{Lem:DR_alteration_invar}
\source{pixtons-formula-abel-jacobi-picard-stack}
We have %$\phi_{\ca L}^*{Pixton}^{g, r=0} = \phi_{\ca L}^*\DRop$. 
$\DRop_{g,\emptyset}(\phi_{\ca L})=
\P_{g,\emptyset,0}^g(\phi_{\ca L}) :\Chow_*(B) \to \Chow_{*-g}(B)$. 
\end{lemma}
\begin{proof}
We apply Lemma \ref{ALTt} to the family $C/B$ to obtain
$$h\colon B' \to B\, , \ \ \ \ C' \to C_{B'}\, .$$ 
Let $\ca L'$ be the pullback of $\ca L$ to $C'$.
After applying \cref{lem:many_markings} with $A = \bd 0 \in \mathbb{Z}^m$, we obtain
\begin{equation}\label{eq:an_equality}
\source{pixtons-formula-abel-jacobi-picard-stack}
\DRop_{g,\bd 0}(\phi_{\ca L'})=\P_{g,\bd 0,d}^g(\phi_{\ca L'}):
\Chow_*(B') \to \Chow_{*-g}(B')\, .
\end{equation}


Since $h$ is proper and surjective, for any $\alpha \in \Chow_*(B)$ there exists $\alpha'\in \Chow_*(B')$ satisfying $h_*\alpha'=\alpha$. If any operational class maps $\alpha'$ to $0$, then it maps $\alpha$ to $0$ because the operation 
commutes with $h_*$. 
% It therefore suffices to prove  
% \begin{equation}
% \left( \DRop_{g,\emptyset}-\P_{g,\emptyset,0}^g \right) (\phi_\ca L) \circ h_*
% \end{equation}
% is the zero map on $\Chow_*(B')$.

% \jocomment{By the compatibilities of operational classes we have
% \begin{equation*}
% \left( \DRop_{g,\emptyset}-\P_{g,\emptyset,0}^g \right) (\phi_\ca L) \circ h_* = h_* \left( \DRop_{g,\emptyset}-\P_{g,\emptyset,0}^g \right) (\phi_\ca L \circ h)
% \end{equation*}
% and the proof below will in fact show $\left( \DRop_{g,\emptyset}-\P_{g,\emptyset,0}^g \right) (\phi_\ca L \circ h)$ since Invariance \invnrIV{} is a statement on equality of operational Chow classes on $B'$. I propose not to drag along these $h_*$ below this point.
% }


% By \eqref{eq:an_equality}, we need only show 
% \begin{equation}\label{eq:pullback_Div}
% \DRop_{g, \emptyset}(\phi_\ca L)\circ h_* = h_* \circ \DRop_{g,\bd 0}(\phi_{\ca L'})
% :\Chow_*(B') \to \Chow_{*-g}(B),
% \end{equation}
% \begin{equation}\label{eq:pullback_pixton}
% \P_{g,\emptyset, 0}^g(\phi_\ca L)\circ h_* = h_* \circ \P_{g,\bd 0, 0}^g(\phi_{\ca L'})
% :\Chow_*(B') \to \Chow_{*-g}(B)
% \, .
% \end{equation}
% First note that for the forgetful map $F : \Picabs_{g,n,0} \to \Picabs_{g,0,0}$ of the markings, we see
% \[\DRop_{g,\bd 0}(\phi_{\ca L'}) = F^*\DRop_{g,\emptyset}(\phi_{\ca L'}) = \DRop_{g,\emptyset}(F \circ \phi_{\ca L'})  \]
% by \cref{Lem:DR_forget_pullback} and similarly
% \[
% \P_{g,\bd 0, 0}^g(\phi_{\ca L'}) = F^*\P_{g,\emptyset, 0}^g(\phi_{\ca L'}) = \P_{g,\emptyset, 0}^g(F \circ \phi_{\ca L'})
% \]
% by \cref{wwwrrr}. 



% \Rcomment{Rewrite without using the forward reference to Sec 7!}
% Second, we use Invariance \invnrIV{}. Indeed, in the notation of \cref{sec:proof_inv_IV} we define a map $\psi\colon B' \to \frak N$ using the data $C' \to C\times_B B'$, $\ca L$. Then $F \circ \phi_{\ca L'} = \sourcemap \circ \psi$ and $\phi_\ca L \circ h = \targetmap \circ \psi$ (where $\targetmap$ and $\sourcemap$ are the maps defined in \cref{sec:proof_inv_IV}), so the equalities  \eqref{eq:pullback_pixton} and \eqref{eq:pullback_Div} follow from \cref{lem:N_invariance_P,lem:N_invariance_DR} respectively.

% \jocomment{To avoid forward references, we need to add to the statement of Invariance \invnrIV{} in \cref{Sect:introinvariance} explicitly the statement of invariance for Pixton's formula (the logic of our proof here requires independently the invariance of $\DRop$ and $\P$). Below is my proposal for the entire end of the proof starting with (75).
% }

It therefore suffices to prove  
\begin{equation}
\left( \DRop_{g,\emptyset}-\P_{g,\emptyset,0}^g \right) (\phi_\ca L) \circ h_*
\end{equation}
is the zero map on $\Chow_*(B')$. By the compatibilities of operational classes we have
\begin{equation*}
\left( \DRop_{g,\emptyset}-\P_{g,\emptyset,0}^g \right) (\phi_\ca L) \circ h_* = h_* \left( \DRop_{g,\emptyset}-\P_{g,\emptyset,0}^g \right) (\phi_\ca L \circ h)
\end{equation*}
and the proof below will in fact show $\left( \DRop_{g,\emptyset}-\P_{g,\emptyset,0}^g \right) (\phi_\ca L \circ h)=0$. 



By \eqref{eq:an_equality}, we need only show 
\begin{equation}\label{eq:pullback_Div}
\source{pixtons-formula-abel-jacobi-picard-stack}
\DRop_{g, \emptyset}(\phi_\ca L\circ h)  =  \DRop_{g,\bd 0}(\phi_{\ca L'})
:\Chow_*(B') \to \Chow_{*-g}(B'),
\end{equation}
\begin{equation}\label{eq:pullback_pixton}
\source{pixtons-formula-abel-jacobi-picard-stack}
\P_{g,\emptyset, 0}^g(\phi_\ca L \circ h) = \P_{g,\bd 0, 0}^g(\phi_{\ca L'})
:\Chow_*(B') \to \Chow_{*-g}(B')
\, .
\end{equation}
For the map $F : \Picabs_{g,n,0} \to \Picabs_{g,0,0}$ forgetting the markings, Invariance \invnrn{} from \cref{Sect:introinvariance} for the double ramification cycle and the Pixton formula shows that
we have
\begin{align*}
\DRop_{g,\bd 0}(\phi_{\ca L'})
%= F^*\DRop_{g,\emptyset}(\phi_{\ca L'}) 
&= \DRop_{g,\emptyset}(F \circ \phi_{\ca L'})  \\
\P_{g,\bd 0, 0}^g(\phi_{\ca L'}) 
% = F^*\P_{g,\emptyset, 0}^g(\phi_{\ca L'}) 
&= \P_{g,\emptyset, 0}^g(F \circ \phi_{\ca L'})
\end{align*}
So we are reduced to showing
\begin{equation} \label{eqn:InvIVapplication}
\source{pixtons-formula-abel-jacobi-picard-stack}
    \DRop_{g, \emptyset}(\phi_\ca L\circ h) = \DRop_{g, \emptyset}(F \circ \phi_{\ca L'})\ \ \text{ and }\ \ \P_{g,\emptyset, 0}^g(\phi_\ca L\circ h) = \P_{g,\emptyset, 0}^g(F \circ \phi_{\ca L'})\ .
\end{equation}


The claims \eqref{eqn:InvIVapplication} follow from
 Invariance \invnrIV{} of \cref{Sect:introinvariance}. As before, let
 $C_{B'}$ be the pullback of $C$ under $h$,  and let
 $\ca L_{B'}$ be the pullback of $\ca L$ to $C_{B'}$. The map 
 $$\phi_\ca L\circ h  : B' \to \Picabs_{g,0,0}$$
 is induced by the data
\[C_{B'} \to B'\, ,\ \ \ \ca L_{B'} \to C_{B'}\, ,\]
whereas $F \circ \phi_{\ca L'} : B' \to \Picabs_{g,0,0}$ is induced by
\[C' \to B'\, ,\ \ \  \ca L'\to C'\,.\]
By construction, we have a partial destabilization $C' \to C_{B'}$ over $B'$, and the line bundle $\ca L'$ is the pullback of $\ca L_{B'}$ under this map. Hence the equalities \eqref{eqn:InvIVapplication} follow from Invariance \invnrIV{} of \cref{Sect:introinvariance}.
\end{proof}
